\clearpage
\phantomsection\addcontentsline{toc}{section}{References}
\begin{thebibliography}{99}\small\raggedright
\bibitem{PaperI} UNNS Research Collective, \emph{Arithmetic Transport and Dimensional Obstructions in UNNS Occurrence Geometry}, Manuscript~I of the programme GEO-HD-1, version~0.4 (2~Nov.\ 2026); file \texttt{PAPER\_I\_TRANSPORT\_v0.4.pdf}, SHA-256 \texttt{19fc10f546696909eef39c75ebdb47b5468d93173f9bd9e82567bd5f60f9eef9}.
\bibitem{PaperII} UNNS Research Collective, \emph{Generative Histories, Arithmetic Recombination, and Higher-Dimensional Route Topology}, Manuscript~II of the programme GEO-HD-1, version~0.4 (2~Nov.\ 2026); file \texttt{PAPER\_II\_HISTORIES\_v0.4.pdf}, SHA-256 \texttt{df34fa94fed9fb33b82c87726ca9e725307723f4e3cb67e1af2f55d7b56801e4}.
\bibitem{PaperIII} UNNS Research Collective, \emph{Geometric Evolution under Collective UNNS Dynamics: Activation, Events, and Topological Prediction}, Manuscript~III of the programme GEO-HD-1, version~0.4 (3~Nov.\ 2026); file \texttt{PAPER\_III\_EVOLUTION\_v0.4.pdf}, SHA-256 \texttt{47fae77fd9c11f0cc7e81447fd697cdbae413e2646564b1d966836ebad14cb8e}.
\bibitem{FoundI} UNNS Substrate programme, \emph{Mathematical Foundations of the UNNS Substrate, Volume~I: Generative Arithmetic and Structural Reachability} (the grammar SELECT-II and the occurrence set). Source: \texttt{ukbbi/UNNS}, commit \texttt{9a021346b2b3d964523440d2737b9d0c4ddbb3e5}, \texttt{projects/UNNS\_MATH\_FOUNDATIONS/Manuscript}.
\bibitem{FoundII} UNNS Substrate programme, \emph{Mathematical Foundations of the UNNS Substrate, Volume~II: Collective Dynamics, Prime Rigidity, and Effective Predictability} (the collective map FIELD-I). Same commit and path as~\cite{FoundI}.
\bibitem{GEOI} GEO-I, Papers~A and~B of the geometry programme (occurrence geometry, region-Graver sets and the realised complexes), cited at commit \texttt{9a021346b2b3d964523440d2737b9d0c4ddbb3e5} as in~\cite{FoundI}.
\bibitem{RefColl} Foundations reference collection v001 (SELECT-I/II, FIELD-I v001--v012; frozen), \texttt{ukbbi/UNNS}, commit \texttt{bb8b12d753c600c597feb00df17fe139c7cc65fa} (23 Oct.\ 2026), \texttt{projects/UNNS\_MATH\_FOUNDATIONS/sources}.
\bibitem{HD1} GEO-HD-1 research collection, snapshot \texttt{a88f31487f3b05c9db297dcd1192d56eea2b0c7d} (the pin named in the programme specification; not independently verified in this work). Frozen releases from which the inherited records are taken, with SHA-256 as computed on the files supplied for Paper~III:
\begin{itemize}[leftmargin=1.4em,itemsep=1pt,topsep=2pt]\footnotesize
\item\texttt{GEO\_HD\_1D\_v001} \texttt{c532b7410f3bf9d437f9f172b3e56029cd9f7bab5b2ac0c0e533713155eaf7a4}; \texttt{GEO\_HD\_1D\_v002} \texttt{8baf463c07bb5285dd53b355aee05a0f841c068ba3ae4794f3702adff6a4bae0}; \texttt{GEO\_HD\_1D\_v003} \texttt{65dbb281e4c75d625b6cf3e4bcb8262d7d42e68d0c476fbba225fcde744a4706}; \texttt{GEO\_HD\_1D\_v004} \texttt{e0390e2a2f88eaeb11042de09fd056741ed98f94318de708d308faa090f05b5b};
\item\texttt{GEO\_HD\_BRIDGE\_v001} \texttt{e99472809928f4a23fb292f60759965159c7a4db5f018d56748ccac328719395}; \texttt{v002} \texttt{012ead97e5b8854cd5218417d46f0e388d927b6d59d5f6e75a8cc343f5dd4955} (with errata for B2-T9 and B2-T7); \texttt{v003} \texttt{6fba87af9f3096772e43f4ac77fa8d55aea9e242eeef241081a987b567d5d5a8}; \texttt{v004} \texttt{cfb31157a39666d1a4bbb33b280ec626c7de2c141e8232da43c8b934b8affff4}; \texttt{v005} \texttt{50e342821d7a5f760c5127df348e97f3a4d4d3bf34bf60e1019d6336f28c5d03}.
\end{itemize}
\bibitem{Graver1975} J.~E. Graver, On the foundations of linear and integer linear programming~I, \emph{Math.\ Programming} \textbf{9} (1975), 207--226.
\bibitem{Sturmfels1996} B.~Sturmfels, \emph{Gr\"obner Bases and Convex Polytopes}, University Lecture Series~8, American Mathematical Society, 1996 (Graver bases, \S4).
\bibitem{Cassels1959} J.~W.~S. Cassels, \emph{An Introduction to the Geometry of Numbers}, Grundlehren der mathematischen Wissenschaften~99, Springer, 1959 (lattices; sublattices as kernels).
\bibitem{DiekertRozenberg1995} V.~Diekert, G.~Rozenberg (eds.), \emph{The Book of Traces}, World Scientific, 1995 (partially commutative monoids).
\bibitem{Hatcher2002} A.~Hatcher, \emph{Algebraic Topology}, Cambridge University Press, 2002 (cellular homology, \S2.2; long exact sequence of a pair, \S2.1).
\bibitem{CarlssonDeSilva2010} G.~Carlsson, V.~de~Silva, Zigzag persistence, \emph{Found.\ Comput.\ Math.} \textbf{10} (2010), 367--405 (zig-zag diagrams of complexes; cited for the shape $K\leftarrow K^\circ\to K'$ only).
\end{thebibliography}
\paragraph{On the references.} Papers~I--III are cited at version~0.4 by theorem number; every such number is generated from the corresponding \texttt{.aux} file (\texttt{scripts/gen\_xref.py}), so a citation that does not resolve is reported at build time. UNNS foundation sources are cited at the commits and archive hashes fixed by the programme; repository access during this work was read-only and the pinned repositories were not fetched. External literature is cited only for the standard notions named and for resemblance of structures; no statement is proved by appeal to an equivalence with those theories.
